% Generated draft: requires mathematical and source review.
\documentclass[11pt]{article}
\usepackage[margin=25mm]{geometry}
\usepackage{fontspec}
\setmainfont{Latin Modern Roman}
\usepackage{amsmath,amssymb,mathtools,unicode-math}
\setmathfont{Latin Modern Math}
\usepackage{enumitem,needspace,xurl,hyperref}
\setlength{\parindent}{0pt}
\setlength{\parskip}{5pt}
\setlength{\emergencystretch}{3em}
\newcommand{\N}{\mathbb{N}}
\newcommand{\Z}{\mathbb{Z}}
\newcommand{\Q}{\mathbb{Q}}
\newcommand{\R}{\mathbb{R}}
\newcommand{\C}{\mathbb{C}}
\newcommand{\F}{\mathbb{F}}
\newcommand{\A}{\mathbb{A}}
\newcommand{\PP}{\mathbb{P}}
\newcommand{\E}{\mathbb{E}}
\newcommand{\eps}{\varepsilon}
\newcommand{\dd}{\,\mathrm{d}}
\begin{document}
\section*{UnsolvedMath ID: 2302062}
\textbf{Source problem code:} catalogue entry
\subsection{Definitions and mathematical statement}
Let $f$ be a rational map of the Riemann sphere $\widehat\C$ of degree at least $2$, or a transcendental entire function on $\C$. Define $f^{\circ1}=f$ and $f^{\circ(n+1)}=f\circ f^{\circ n}$. Its Fatou set $\mathcal F(f)$ is the open set on which the iterates form a normal family, using the spherical metric for meromorphic maps. A Fatou component is a connected component of $\mathcal F(f)$. The source denotes this normality set by $C$ and notes that it is proper, open, and forward invariant.

The supplied definition calls a component $U$ wandering when
\[
f^{\circ k}(U)\cap f^{\circ n}(U)=\varnothing
\quad\text{for all distinct integers }k,n\geq1.
\]
Jakobson's question is whether there exists a rational map of degree at least $2$ with such a wandering domain. The degree-$0$ and degree-$1$ rational cases are explicitly excluded. The source separately reports Baker's construction of a transcendental entire function with wandering domains, so the rational and transcendental cases must not be conflated.

The supplied literature resolution is Sullivan's no-wandering-domains theorem: every Fatou component of a rational map of degree at least $2$ is eventually periodic. More precisely, if $U_j$ denotes the component containing $f^{\circ j}(U)$, there are integers $q\geq0$, $p\geq1$ with $U_{q+p}=U_q$. Consequently the rational-map existence question has a reported negative answer. This theorem is recorded as supplied literature, not proved here.
\subsection{Short English statement}
The question asks whether a rational map can have a Fatou component whose successive images wander without repetition. The supplied record cites Sullivan's theorem giving a negative answer for rational maps of degree at least $2$, while retaining Baker's contrasting entire-function example.
\subsection{Sources}
\begin{itemize}
\item Local UnsolvedMath registry, permanent ID 2302062, exact source code catalogue entry, ``Research Problems in Function Theory --- Problem 2.62''; status solved, classification SOLVED-IN-LITERATURE.
\item Hayman, \emph{Research Problems in Function Theory} (2018), Problem 2.62, attributed in the item to I. N. Baker and reporting Jakobson's question, \url{https://arxiv.org/abs/1809.07200}. The record identifies extraction from public arXiv source TeX, accessed July 29, 2026; the 2018 source update was retained for status triage. The statement credits Baker with a transcendental entire example, without a separate bibliographic citation.
\item D. Sullivan, ``Quasiconformal homeomorphisms and dynamics I: solution of the Fatou--Julia problem on wandering domains,'' \emph{Annals of Mathematics} 122 (1985), 401--418, cited in the supplied research summary for eventual periodicity of every Fatou component of rational maps of degree at least $2$.
\end{itemize}
\end{document}
